% s-026-p026-m001; original PDF 26, 27, 28
\Aphorism{s-026-p026-m001}{17}
Die Stimmung, in der man zum Kurpfuscher oder Magnetiseur geht, nachdem man mit den wissenschaftlichen Ärzten „durch“ ist: etwas Ähnliches stelle ich mir als Grundton jenes Gefühlsaccordes von Neugier, Hoffnung, Verachtung, Widerwillen vor, mit dem das ausklingende \Foreign{imperium Romanum} ins Christenthum hinübermodulirte. Die neurasthenischen und conträr sexualen Römer der Kaiserzeit, die alle physiologischen Voraussetzungen der Weltmüdigkeit allzusehr aus der Nähe kannten — wie mussten sie die skeptischen Ohren spitzen bei der christlichen Deutung jener Zustände, einer Deutung, die nicht ohne erhabene Reclame in die Empfehlung einer Wunderkur auslief! Solche Ungeheuer von dämonischer Wühlkraft und feinspüriger Selbstquälerei wie Nero und Julian (gegen die ein Caligula freilich nur wie ein blutrünstiger Fleischer erscheint), Menschen, die lange vor Schopenhauer das Mitleid mit allem Lebenden und den untilgbaren Leidensfonds der Menschheit kannten — wie musste sie der Anblick jener theoretischen Naiven würgen und kitzeln, die mit seiltänzerischer Unbedenklichkeit über den Abgründen spazieren gingen und von ihren beiden \Foreign{civitates} so sicher und unbedingt sprachen wie der Mensch von rechts und links! Das war ein neuer Reiz für abgewirthschaftete Nerven, ein neues Denkproblem für die todtgehetzte Logik unverbesserlicher Staatsfanatiker und „Weltkinder“; hier war jenem Ruhebedürfniss eines späten Zeitalters, dessen melancholischen Ausdruck Horaz gefunden hat, ein letztes erlöstes Aufathmen versprochen, aber um den Preis einer radicalen Umkehr und Selbstvergewaltigung. Jede Praktik und Gebärde der neuen Secte widersprach römischer Tradition — auch die Quacksalber und Naturärzte thun es nie unter einer völligen „Regeneration“ und Abtödtung des alten Adam. Dem römischen Recht wurde zugemuthet, ein Processverfahren ernst zu nehmen, wo sich das Leben selbst vor einigen Lebenden zu verantworten hatte, d.h. wo die Richter selbst Partei waren. Der römischen Baupolizei und Architectur zeigte man jene Polizeiwidrigkeit von Bauwerk, von dem aus die Erde, auf der es selbst stand, aus den Angeln gehoben und abgewogen werden sollte: zu dieser physikalischen Absurdität wollte man, ein Viertel Jahrtausend nach Archimedes, ein realistisches und constructives Volk überreden. Man hatte damit den Erfolg, den heute magnetische Curen und Geheimmittel haben; kein Mensch glaubt daran, aber man versucht sie, als das Letzte oder Vorletzte. Nur war man in Rom vorsichtig genug, das Geheimmittel zuerst seinem Erfinder einzugeben, sowie Phalaris den ehernen Stier sogleich an seinem Verfertiger Perilaos erprobte. Die bizarren und schauerlichen Einzelheiten, mit denen sich der römische Geschmack jede Christenverfolgung zurechtwürzte, sind eine possenhafte \Foreign{deductio ad absurdum}, eine Art Rechenprobe auf die vorausgeahnte Unmöglichkeit und Unwirksamkeit der christlichen Neuerung; man wollte den neu aufgetauchten Glauben, von dessen Allmacht Wunderdinge gerühmt wurden, auf einer Schneide schweben sehen, wo er sich entscheiden \Sperr{musste}! „Schafft euch ab, wenn ihr könnt: hier habt ihr Grund genug dazu! Setzt eure gepredigte Widernatur in Thaten um, wenn ihr könnt: wir selber halten es in unserer Natur nicht mehr aus und wären die Ersten, euch zu folgen! Zeigt, ob euer Narcoticum über Rost und Flamme vorhält, vielleicht verfängt es dann auch bei unseren Nothständen!“ Dergleichen muss Nero empfunden haben, als er die lebenden Fackeln entzündete; das ist die heimliche Logik jener grausamen Experimentirsucht, mit der die römische \Foreign{décadence} über die ersten christlichen Organisationen herfiel.
\par

